\chapter{Introdução}
\addcontentsline{toc}{chapter}{Introdução}


A comunicação dos bancos centrais desempenha papel fundamental na formação das 
expectativas dos agentes econômicos e constitui parte relevante da própria 
transmissão da política monetária, uma vez que informações sobre perspectivas 
econômicas e possíveis trajetórias futuras dos juros podem influenciar as 
expectativas dos participantes do mercado \citep{woodford2005central}.
No caso do Federal Reserve (Fed), autoridade monetária dos Estados Unidos, 
essa comunicação ocorre por meio de diferentes documentos, como comunicados e 
atas das reuniões do Federal Open Market Committee (FOMC), discursos de seus 
dirigentes e outros materiais institucionais. Essas divulgações são amplamente 
acompanhadas por participantes do mercado financeiro e 
pesquisadores, que buscam identificar sinais sobre a orientação futura da política 
monetária e antecipar possíveis decisões relacionadas à taxa básica de 
juros \citep{blinder2008central}. Evidências empíricas indicam que os mercados 
reagem não apenas às decisões correntes sobre a taxa de juros, mas também às 
informações contidas nos comunicados 
do FOMC acerca da trajetória futura da política monetária \citep{gurkaynak2005actions}.

Apesar da relevância dessas comunicações, a interpretação de seu conteúdo ainda 
representa um desafio significativo. O volume elevado de documentos, aliado à 
complexidade da linguagem utilizada, dificulta a extração sistemática de informações e 
a identificação de padrões consistentes ao longo do tempo. Estudos indicam que análises 
tradicionais baseadas em leitura manual ou métodos simplificados apresentam limitações 
na captura de nuances semânticas e relações implícitas nos textos \citep{loughran2011liability}.
Dessa forma, surge a problemática de como analisar, de maneira estruturada e automatizada, 
os conteúdos textuais produzidos pelo Federal Reserve, de modo a gerar insights úteis para a 
compreensão e antecipação da política monetária.

Essa problemática insere-se em um contexto onde documentos textuais são tratados como fontes 
de informação quantitativa para a investigação de fenômenos econômicos \citep{gentzkow2019text}. 
Nesse contexto, os avanços recentes em técnicas de processamento de linguagem natural 
(Natural Language Processing - NLP), especialmente os chamados modelos de linguagem de 
grande porte (Large Language Models - LLMs), ampliaram significativamente as possibilidades 
de análise automatizada de textos complexos. Esses modelos são capazes de identificar padrões 
semânticos, extrair informações e sintetizar conteúdos presentes em grandes volumes de 
documentos textuais \citep{brown2020language}.

De forma mais específica \cite{hansen2024can} investiga a capacidade 
de modelos GPT em interpretar a linguagem técnica utilizada pelo Federal 
Reserve em suas comunicações de política monetária. Os autores avaliam o desempenho 
desses modelos na classificação da orientação de comunicados do FOMC em comparação 
com avaliações humanas, indicando que os modelos GPT apresentam melhora relevante 
em relação a métodos tradicionais de classificação textual. Esse resultado reforça 
o potencial dos LLMs para capturar nuances semânticas presentes na comunicação de 
bancos centrais, especialmente na identificação de sinais associados a posturas 
mais expansionistas ou restritivas de política monetária.

\cite{silva2025text} utilizam modelos de linguagem de grande porte para analisar 
comunicações de diferentes bancos centrais, classificando os documentos segundo 
dimensões como tópico, orientação da política monetária e sentimento. 
A partir dessas classificações, os autores constroem medidas quantitativas capazes 
de sintetizar a orientação presente nas comunicações e investigam sua relação com 
mudanças futuras na política monetária. Os resultados demonstram o potencial dos 
LLMs não apenas para interpretar grandes volumes de documentos textuais, mas também 
para transformar informações qualitativas presentes na comunicação dos bancos centrais 
em indicadores quantitativos passíveis de análise ao longo do tempo.

Paralelamente, técnicas baseadas em grafos permitem representar e analisar relações 
entre informações que não são facilmente observadas quando os elementos são considerados 
de forma isolada. Em uma rede, os nós representam os elementos de interesse e as 
conexões descrevem as relações existentes entre eles, possibilitando a investigação de 
propriedades como proximidade, centralidade e formação de comunidades 
\citep{needham2019graph}. 

No contexto deste estudo, os sinais de política 
monetária identificados nas comunicações do Federal Reserve são representados em uma estrutura 
de grafo a partir de suas relações de similaridade semântica. Essa representação permite 
investigar não apenas a orientação individual de cada sinal, mas também sua posição 
relativa na rede, sua associação com outros sinais semanticamente semelhantes e a 
formação de agrupamentos que podem refletir padrões recorrentes na comunicação da autoridade 
monetária.
 
Paralelamente, técnicas de análise de grafos permitem representar e explorar 
relações entre informações extraídas de documentos. Algoritmos de grafos podem 
identificar proximidades e agrupamentos, 
revelando padrões difíceis de observar quando os dados são analisados 
isoladamente \citep{needham2019graph}. No contexto deste estudo, relações de 
similaridade semântica entre os sinais monetários extraídos das comunicações do 
Federal Reserve são utilizadas para construir uma rede sobre a qual podem ser 
analisadas estruturas de comunidades e medidas de centralidade. Essas propriedades 
estruturais servem para complementar a classificação atribuída 
individualmente aos sinais, uma vez que consideram também o contexto relacional 
em que cada sinal está inserido. 

Embora estudos recentes tenham demonstrado o potencial dos modelos de 
linguagem para interpretar e quantificar a comunicação de bancos centrais, 
permanece menos explorada a utilização conjunta dessas técnicas com métodos de 
análise de grafos para investigar as relações existentes entre os sinais 
monetários extraídos dos documentos. Dessa forma, este estudo busca investigar 
se a combinação entre modelos de linguagem de grande porte 
e técnicas de análise de grafos pode contribuir para a construção de um indicador 
experimental da orientação da política monetária do Federal Reserve, avaliando sua 
capacidade de sintetizar padrões presentes nas comunicações e sua relação com 
decisões posteriores do FOMC.